%% LyX 2.5.1 created this file.  For more info, see https://www.lyx.org/.
%% Do not edit unless you really know what you are doing.
\documentclass[english,hebrew]{article}
\usepackage[HE8,T1]{fontenc}
\usepackage[utf8]{inputenc}
\usepackage{babel}
\usepackage[unicode=true]
 {hyperref}
\begin{document}
\title{מבוא לטופולוגיה קבוצתית: משפט טיכונוף}
\maketitle
\begin{description}
\item [{קטגוריות:}] טופולוגיה
\item [{תגים:}] משפט טיכונוף
\item [{מזהה:}] \L{tychonoff\_theorem}
\end{description}
משפט טיכונוף הוא מהמשפטים הללו שקלים לניסוח אבל טריקיים להוכחה. הוא
אומר ש\textbf{מכפלה של מרחבים קומפקטיים היא קומפקטית}, כשכאן \textquotedblleft מרחבים\textquotedblright{}
הם מרחבים טופולוגיים ו\textquotedblright מכפלה\textquotedblright{}
היא מרחב המכפלה \textbf{עם טופולוגיית המכפלה} )עם טופולוגיית הקופסה
המשפט לא נכון(. זה משפט חזק מאוד בגלל שהמכפלה יכולה להיות של קבוצה
כלשהי של מרחבים, אפילו קבוצה לא בת מניה שלהם; אבל זו גם הסיבה שבהוכחה
אנחנו נזקקים לכלי חזק למדי - \textbf{אקסיומת הבחירה} )לא נשתמש בה
ישירות אלא בלמה של צורן ששקולה לה(. אפשר גם לתת הוכחה פשוטה יותר ונטולת
אקסיומת הבחירה לכך שמכפלה \textbf{סופית} של מרחבים קומפקטיים היא קומפקטית
אבל אני אוותר על זה כאן.

ההוכחה יכולה להיות טיפה מאתגרת כשנתקלים בה לראשונה, אבל היא די דומה
באופיה להוכחות דומות שמתבססות על בניה כמו של הלמה של צורן. הראיתי
בבלוג את \href{https://gadial.net/2013/02/26/godel_completeness_proof_2/}{משפט השלמות של גדל}
שבו יש משהו כזה, וגם \href{https://gadial.net/2013/06/21/ultrafilters_and_ultraproducts/}{הראיתי הוכחה}
לקיום על-מסנן שגם היא מאוד דומה. בפוסט הקודם הזכרתי את איך שמשפט טיכונוף
יכול להוכיח את \textbf{משפט הקומפקטיות} בלוגיקה, ואותו משפט יכול לנבוע
גם ממשפט השלמות של גדל וגם מקיום על-מסנן. במילים אחרות - ראינו כבר
את החדק והזנב של הסיפור הזה, היום אנחנו הולכים לראות את האוזניים.

יהיה נוח להשתמש בהגדרה של קומפקטיות בעזרת חיתוכים, אז בואו נזכיר אותה:
מרחב \L{$X$} הוא קומפקטי אם לכל אוסף של קבוצות סגורות שמקיים את \textbf{תכונת
החיתוך הסופי} יש חיתוך לא ריק. 

כלומר, פורמלית: לכל אוסף \L{$\left\{ C_{\alpha}\right\} _{\alpha\in\Lambda}$}
של קבוצות סגורות ב-\L{$X$} כך שלכל תת-אוסף \textbf{סופי} \L{$\Lambda^{\prime}\subseteq\Lambda$}
מתקיים \L{$\bigcap_{\alpha\in\Lambda^{\prime}}C_{\alpha}\ne\emptyset$}
מתקיים ש-\L{$\bigcap_{\alpha\in\Lambda}C_{\alpha}\ne\emptyset$}.

הרעיון \textquotedblleft בגדול\textquotedblright{} של ההוכחה הוא מאוד
פשוט; הוא הולך להיכשל בשתי צורות שונות, אחת לא כזו נוראית ואחת יותר
מורכבת, אבל הוא מאוד פשוט. נניח לצורך הפשטות ש-\L{$X=A\times B$}
ו-\L{$A,B$} קומפקטיים. נסתכל על אוסף \L{$\left\{ C_{\alpha}\right\} _{\alpha\in\Lambda}$}
של קבוצות סגורות ב-\L{$X$} שמקיים את תכונת החיתוך הסופי. אנחנו רוצים
למצוא איבר \L{$\left(a,b\right)\in A\times B$} כך ש-\L{$\left(a,b\right)\in\bigcap_{\alpha\in\Lambda}C_{\alpha}$}.
הנה דרך פשוטה )וכאמור, שגויה כל עוד מבצעים את זה נאיבית( לעשות את
זה:
\begin{itemize}
\item נסתכל על \textbf{ההטלה} של הקבוצות באוסף על \L{$A$}, כלומר על \L{$\pi_{A}\left(C_{\alpha}\right)$}
לכל \L{$\alpha\in\Lambda$}. קיבלנו אוסף של קבוצות ב-\L{$A$} והוא
מקיים את תכונת החיתוך הסופי כי ה-\L{$C_{\alpha}$} המקוריים קיימו
אותה )הטענה הספציפית הזו \textbf{כן נכונה} אבל גם אותה אוכיח בהמשך(.
לכן קיים \L{$a\in\bigcap_{\alpha\in\Lambda}\pi_{A}\left(C_{\alpha}\right)$}.
\item נעשה את אותו דבר עם \L{$B$} ונמצא \L{$b\in\bigcap_{\alpha\in\Lambda}\pi_{B}\left(C_{\alpha}\right)$}.
\item קיבלנו איבר \L{$\left(a,b\right)\in A\times B$} כך ש-\L{$\left(a,b\right)\in\bigcap_{\alpha\in\Lambda}C_{\alpha}$}
וסיימנו.
\end{itemize}
למה זה לא עובד? ראשית, כי אני משתמש בהגדרת החיתוך של קומפקטיות על
\textbf{ההטלות} \L{$\pi_{A}\left(C_{\alpha}\right)$}, אבל בשביל שאוכל
להשתמש בה, ההטלות צריכות להיות קבוצות \textbf{סגורות}. והן לא. למרבה
המזל זו לא בעיה רצינית כי אפשר יהיה להסתכל על הסגור שלהן, \L{$\overline{\pi_{A}\left(C_{\alpha}\right)}$},
והוא כן יהיה קבוצה סגורה ונקבל אוסף של קבוצות סגורות שמקיים את תכונת
החיתוך הסופי, כך שנקבל את ה-\L{$a$} שלנו בסוף.

אבל הבעיה המהותית הרבה יותר היא זו - אוקיי, מצאנו \L{$a$} ששייך לכל
ההטלות של הקבוצות על \L{$A$}; מצאנו \L{$b$} ששייך לכל ההטלות של
הקבוצות על \L{$B$}; אבל למה זה מבטיח ש-\L{$\left(a,b\right)$} שייך
לכל הקבוצות המקוריות? וזו בעיה משמעותית כי \textbf{הוא לא}. בואו נראה
דוגמא קונקרטית פשוטה לזה. נניח ש-\L{$X=\left[0,1\right]^{2}$}, כלומר
\L{$A=B=\left[0,1\right]$}, וניקח קבוצה סגורה \textbf{אחת}, האלכסון
\L{$C=\left\{ \left(x,x\right)\ |\ x\in\left[0,1\right]\right\} $}.
ההטלה של \L{$C$} על \L{$A$} היא פשוט \L{$A$} כולה ובדומה גם ההטלה
של \L{$C$} על \L{$B$}. עכשיו נניח שנבחר \L{$a=0$} ו-\L{$b=1$},
נקבל את הנקודה \L{$\left(0,1\right)$} שכמובן לא על האלכסון. זה מחרב
לנו את כל הרעיון. מחרב בצורה כזו של \textquotedblleft אוקיי, למה בעצם
חשבנו שהכיוון הזה יכול לעבוד, חזרה לשולחן השרטוט\textquotedblright .

אבל לא! אולי אפשר עוד קצת להתעקש. הרי אחרי שבחרתי \L{$a=0$} זה לא
שהסיפור היה אבוד, הייתי יכול לבחר \L{$b=0$} ולקבל נקודה שבאמת שייכת
ל-\L{$C$}. אז אולי הייתי צריך שהבחירה \L{$a=0$} \textbf{תגביל} איכשהו
את הבחירה שלי ב-\L{$B$}? אולי אפשר לעשות את זה, אבל כשיש לנו מספר
\textbf{אינסופי} של מרחבים שאנחנו מכפילים זה בזה, נצטרך סוג של אינדוקציה
כדי לפרמל את הטיעון; וכשיש לנו מספר \textbf{לא בן מניה} של מרחבים
נצטרך להשתמש במה שנקרא \textbf{אינדוקציה טרנספיניטית} )או בשמה העברי,
\textquotedblleft על-סופית\textquotedblright ; אני מדבר עליה \href{https://gadial.net/2023/01/18/classes_and_transfinite/}{כאן}(.
הכיוון הזה בהחלט \textbf{יכול} לעבוד, אבל אפשר לטפל בבעיה יותר באלגנטיות
- למעשה, זה מה שהופך את ההוכחה הזו לאחת מהחביבות עלי במתמטיקה. הרעיון
של \textquotedblleft בואו נגביל את מה שאפשר לבחור ב-\L{$B$} כי בבירור
היה לנו \textbf{יותר מדי חופש בחירה} וזה פעל לרעתנו\textquotedblright{}
הוא נכון, אבל אפשר לטפל בו על ידי הגבלה של חופש הבחירה \textbf{מראש},
לפני שמתחילים להסתכל על הטלות.

עכשיו, מה זו \textquotedblleft הגבלה\textquotedblright ? כשבחרתי את
\L{$a$}, ההגבלה הייתה שיתקיים \L{$a\in\bigcap_{\alpha\in\Lambda}\pi_{A}\left(C_{\alpha}\right)$}.
כלומר, הקבוצות \L{$C_{\alpha}$} המקוריות שלי היו אחראיות ליצירת האילוץ
על \L{$a$}. הרעיון במשפט טיכונוף הוא \textbf{להוסיף עוד קבוצות כאלו}
ולכן להחמיר את ההגבלה על ה-\L{$a$} שאפשר למצוא. כמובן שהוספת קבוצות
כאלו לא יכולה להפריע, במובן זה שאם נמצא \L{$a$} ששייך לחיתוך הקבוצות
כולל אלו שהוספנו במהלך ההרחבה, אז הוא יהיה שייך גם לחיתך הקבוצות המקוריות
בלבד.

אוקיי, אז נוסיף עוד קבוצות. אבל כמה? אילו? וכאן יש את תשובת המחץ:
\textbf{נוסיף את המקסימום האפשרי שנוכל} כי בדיוק המקסימום הזה הולך
ליצור את האילוץ האולטימטיבי שיגרום להכל לעבוד; ולגבי ה\textquotedblright אילו\textquotedblright ,
ובכן: \textbf{אין לנו מושג} איזה קבוצות הולכות להתווסף. רק שזה יהיה
אפשרי. הרי לנו לב הבעיה, שתמיד חוזר על עצמו בהוכחות שמסתמכות על הלמה
של צורן/אקסיומת הבחירה: מה שאנחנו עושים הוא לא \textbf{קונסטרוקטיבי}.

בואו ניזכר איך הקטע הזה של הלמה של צורן עובד בכללי. אם יש לנו קבוצה
סדורה חלקית \L{$\left(A,\le\right)$}, הלמה של צורן מראה שבתנאים מסויימים
קיים ל-\L{$A$} \textbf{איבר מקסימלי}, כלומר \L{$M\in A$} כך שאין
\L{$a\in A$} שמקיים \L{$M<a$} )שימו לב שזה לא אותו דבר כמו \L{$a\le M$}
לכל \L{$a\in A$} כי לא כל איבר של \L{$A$} בהכרח יהיה ניתן להשוואה
עם \L{$M$}, יחס הסדר הוא \textbf{חלקי}(. מה התנאים שצריך כדי שדבר
כזה יקרה? ובכן, אפשר לחשוב איך היינו מנסים למצוא איבר מקסימלי: היינו
לוקחים \L{$a\in A$} כלשהו )כלומר, \L{$A\ne\emptyset$}, נתחיל מזה(
ואם הוא מקסימלי, סיימנו. אחרת יש \L{$b\in A$} כך ש-\L{$a<b$}. אם
\L{$b$} מקסימלי, סיימנו, אחרת יש \L{$c$} כך ש-\L{$a<b<c$} וכן הלאה.
מה שקורה הוא שאם אנחנו לא מוצאים איבר מקסימלי, אנחנו מקבלים \textbf{שרשרת}:
תת-קבוצה של \L{$A$} שכל האיברים שלה ניתנים להשוואה זה עם זה )מה שכאמור
- לא בהכרח מתקיים לכל \L{$A$} כי הסדר הוא חלקי(. בצורה הזו אפשר לבנות
שרשרת שהיא סדרה \L{$a_{1},a_{2},a_{3},\ldots$}. אם \L{$A$} הייתה
סופית, זה היה מספיק; מתישהו היו נגמרים לנו האיברים והאחרון היה האיבר
המקסימלי. אבל אם \L{$A$} אינסופית צריך איכשהו להמשיך - לקבל את \textquotedblleft האיבר
הבא\textquotedblright{} שבסוף הסדרה הזו. עבור מי שבקיאים בסודרים, הרעיון
פה הוא לבנות איבר בסדרה לכל סודר קיים, ובשביל הסודרים \textbf{הגבוליים}
צריך איכשהו למצוא איבר שגדול יותר מכל מה שהיה עד כה. כאן נכנסת ההנחה
של צורן: \textbf{לכל שרשרת יש חסם מלעיל}, כלומר לכל שרשרת אנחנו יכולים
למצוא איבר שגדול מכל איבריה. עם ההנחה הזו, אפשר לבנות שרשרת כל כך
ארוכה שמתישהו היא תהיה חייבת למצות את האיברים של \L{$A$} ולכן שנקבל
איבר מקסימלי. כאן גם אקסיומת הבחירה נכנסת לתמונה - בשביל לבנות את
השרשרת הזו צריך \textbf{לבחור} איברים להוסיף וזה מצריך את האקסיומה,
אם כי אני לא אכנס לדקויות של זה הפעם.

בהרבה שימושים, ובפרט בשימוש הנוכחי שלנו, הקבוצה הסדורה היא משפחה של
\textbf{קבוצות} ויחס הסדר הוא פשוט יחס ההכלה הרגיל של קבוצות. אני
אסמן את זה ככה: \L{$\left(\Phi,\subseteq\right)$} כאשר \L{$\Phi$}
היא \textquotedblleft משפחה של קבוצות\textquotedblright{} )שזו דרך
אחרת לומר \textquotedblleft זו קבוצה של קבוצות אבל בואו נגיד שהיא
\textbf{משפחה} כדי לא להתבלבל בינה ובין הקבוצות שהיא מכילה\textquotedblright ,
וזה גם למה אני משתמש באות יוונית גדולה כמו \L{$\Phi$}(. בדרך כלל
יהיה קריטריון כלשהו שקובע שייכות ל-\L{$\Phi$}. עכשיו, בהינתן שרשרת
ב-\L{$\Phi$} הדרך הסטנדרטית לבנות לה חסם מלעיל היא לאחד את כל איבריה;
זה בוודאי ייתן קבוצה חדשה שמכילה את הקודמות, האתגר יהיה להוכיח שגם
היא מקיימת את קריטריון השייכות ל-\L{$\Phi$}.

אצלנו יש רמה נוספת של סיבוך: כל איבר של \L{$\Phi$} שלנו הוא משהו
מהצורה \L{$\left\{ C_{\alpha}\right\} _{\alpha\in\Lambda}$}, שכיניתי
\textbf{אוסף} של קבוצות. נוח לחשוב עליו בתור קבוצה של קבוצות )אבל
אפשר גם לחשוב עליו בתור פונקציה מ-\L{$\Lambda$} לקבוצות( ואז יש לנו
\textbf{ארבע} רמות של איברים:
\begin{itemize}
\item \textbf{איברים} של המרחב הטופולוגי \L{$X$} שלנו, שאני מסמן באותיות
לטיניות קטנות, למשל \L{$a,b\in X$}.
\item \textbf{קבוצות סגורות} של איברים במרחב הטופולוגי \L{$X$}, שאני מסמן
באותיות לטיניות גדולות, למשל \L{$C$}.
\item \textbf{אוספים} של קבוצות סגורות כאלו, שאני אסמן באותיות לטיניות קליגרפיות,
למשל \L{$\mathcal{C}$} כי כמובן שזה לא מבלבל בכלל.
\item \textbf{משפחה} \L{$\Phi$} שכוללת את כל האוספים \L{$\mathcal{C}$}
שמקיימים את תכונת החיתוכים הסופיים.
\end{itemize}
מה שאנחנו צריכים לעשות כדי להשתמש בלמה של צורן הוא להראות שאם יש לנו
שרשרת של קבוצות \L{$\mathcal{C}$} שמקיימות את תכונת החיתוכים הסופיים,
אז גם האיחוד שלהן מקיים את תכונת החיתוכים הסופיים. פורמלית, בואו נסמן
ב-\L{$\Psi$} את השרשרת ונגדיר \L{$\mathcal{D}=\bigcup\Psi$} אז לכל
\L{$\mathcal{C}\in\Psi$} מתקיים \L{$\mathcal{C}\subseteq\mathcal{D}$},
והאתגר שלנו הוא להבין למה \L{$\mathcal{D}$} מקיימת את תכונת החיתוכים
הסופיים )מעבר לאתגר של להבין מה הולך עם כל האותיות הללו(.

מה שאנחנו צריכים להראות הוא שלכל תת-קבוצה \textbf{סופית} של \L{$\mathcal{D}$}
יש חיתוך לא ריק. אני אסמן תת-קבוצה כללית כזו ב-\L{$\left\{ D_{1},D_{2},\ldots,D_{n}\right\} $}.
הסופיות הזו היא בדיוק מה שיאפשר לנו להשתמש בטריק הסטנדרטי שבו משתמשים
בהוכחות כאלו של הלמה של צורן שבהן פשוט מאחדים את אברי השרשרת שלנו.
מכיוון שכל \L{$D\in\mathcal{D}$} הגיע מאחת מהקבוצות שהשתתפו באיחוד
)כלומר, מקבוצה כלשהי מהשרשרת \L{$\Psi$}( אז בואו נסמן ב-\L{$\mathcal{C}_{i}$}
את הקבוצה שממנה \L{$D_{i}$} הגיע, \L{$D_{i}\in\mathcal{C}_{i}$}.
עכשיו, ה-\L{$\mathcal{C}$}-ים הללו הם איברים של \textbf{שרשרת}, כלומר
הם סדורים בסדר מלא ביחס להכלה. אז אפשר לבחור מלכתחילה את האינדקסים
מ-\L{$1$} עד \L{$n$} כך שיתקיים

\L{$\mathcal{C}_{1}\subseteq\mathcal{C}_{2}\subseteq\ldots\subseteq\mathcal{C}_{n}$}

אבל זה אומר \textbf{שכל} האיברים \L{$D_{1},\ldots,D_{n}$} שייכים
כולם אל \L{$\mathcal{C}_{n}$}; ומכיוון ש-\L{$\mathcal{C}_{n}$} מקיימת
את תכונת החיתוכים הסופיים, אז \L{$\bigcap^{n}_{i=1}D_{i}\ne\emptyset$},
סיימנו; הוכחנו שלכל שרשרת יש חסם מלעיל ולכן יש איבר מקסימלי ב-\L{$\Phi$}.

נסכם: הוכחנו שבמרחב הטופולוגי \L{$X$}, לכל אוסף \L{$\mathcal{C}$}
של תת-קבוצות של \L{$X$} שמקיים את \textbf{תכונת החיתוכים הסופיים}
יש קבוצה \L{$\mathcal{C}^{\prime}$}כך ש-\L{$\mathcal{C}\subseteq\mathcal{C}^{\prime}$}
ו-\L{$\mathcal{C}^{\prime}$} היא \textbf{מקסימלית} - כל איבר שנוסיף
לה יוביל לכך שהיא לא תקיים את תכונת החיתוכים הסופיים.

מה שעדיין לא ברור היא איך המקסימליות הזו עוזרת לנו. התשובה לזה פשוטה:
אם \L{$\mathcal{C}$} מקסימלית, יותר קל להראות שדברים שייכים אליה
וזה בדיוק מה שנזדקק לו בהמשך. ליתר דיוק נזדקק לשתי התכונות הבאות:
\begin{itemize}
\item אם \L{$\mathcal{C}$} מקסימלית ביחס לתכונת החיתוכים הסופיים אז היא
סגורה תחת חיתוכים סופיים )כלומר, חיתוך סופי של איברים שלה גם הוא שייך
אליה(.
\item אם \L{$\mathcal{C}$} מקסימלית ביחס לתכונת החיתוכים הסופיים ואם \L{$D$}
היא קבוצה שיש לה חיתוך לא ריק עם \textbf{כל} אבר ב-\L{$\mathcal{C}$},
אז \L{$D\in\mathcal{C}$}.
\end{itemize}
התכונה הראשונה פשוטה מאוד להוכחה. נניח ש-\L{$D=C_{1}\cap\ldots\cap C_{n}$}
הוא חיתוך סופי של אברי \L{$\mathcal{C}$}, ונראה שגם אם מוסיפים את
\L{$D$} אל \L{$\mathcal{C}$} היא עדיין שומרת על תכונת החיתוכים הסופיים
שלה - כלומר \L{$D$} היה חייב מראש להיות איבר שלה, אחרת היינו מקבלים
סתירה למקסימליות )כי במקרה כזה, \L{$\mathcal{C}\cup\left\{ D\right\} $}
הייתה קבוצה שמקיימת את תכונת החיתוכים הסופיים \textbf{וגם} מכילה ממש
את \L{$\mathcal{C}$}(.

למה אחרי שמוסיפים את \L{$D$} תכונת החיתוכים הסופיים משתמרת? כי נניח
שאנחנו מסתכלים על החיתוך הסופי \L{$C^{\prime}_{1}\cap\ldots\cap C^{\prime}_{m}\cap D$}.
אפשר פשוט לכתוב במקום \L{$D$} את הקבוצות שהרכיבו אותו ולקבל \L{$C^{\prime}_{1}\cap\ldots\cap C^{\prime}_{m}\cap C_{1}\cap\ldots\cap C_{n}$}
וזה חיתוך סופי של אברי \L{$\mathcal{C}$} ולכן לא ריק. זה מסיים את
הוכחת הטענה הראשונה.

את הטענה השניה קל להוכיח בעזרת הראשונה. גם פה, הרעיון הוא להראות שלהוסיף
את \L{$D$} אל \L{$\mathcal{C}$} לא פוגע בתכונת החיתוכים הסופיים.
זה לא נראה מובן מאליו: אם \L{$C_{1},C_{2}$} שתי קבוצות שונות אז סתם
להגיד שקיים \L{$y_{1}\in C_{1}\cap D$} וקיים \L{$y_{2}\in C_{2}\cap D$}
לא עוזר לי - אני רוצה להשתכנע בכך ש-\L{$C_{1}\cap C_{2}\cap D$} לא
ריקה, ואם \L{$y_{1}\ne y_{2}$} בהחלט ייתכן שאף אחד מהם לא ייכנס לחיתוך
הזה )כי ייתכן ש-\L{$y_{2}\notin C_{1}$} וההפך(. אלא מה - כמו שראינו
קודם, לא צריך להתעסק עם \textbf{הרבה} קבוצות שונות. באופן כללי אם
יש לנו את החיתוך \L{$D\cap C_{1}\cap\ldots\cap C_{n}$} אז מהטענה
הקודמת, כל \L{$C_{1}\cap\ldots\cap C_{n}$} הוא בעצמו קבוצה ב-\L{$\mathcal{C}$}
ולכן החיתוך שלו עם \L{$D$} לא ריק, וסיימנו.

עכשיו אפשר לעבור להוכחה של משפט טיכונוף עצמו. בואו ניקח מרחב מכפלה
\L{$X=\prod A_{i}$} כך שכל \L{$A_{i}$} הוא קומפקטי )אני לא מניח
כלום על מספר המוכפלים - יכול להיות סופי, בן מניה, לא בן מניה, מה שתרצו(.
ניקח אוסף \L{$\mathcal{C}$} של קבוצות \textbf{סגורות} ב-\L{$X$}
שמקיים את תכונת החיתוכים הסופיים )החיתוך של כל מספר סופי של איברים
של \L{$\mathcal{C}$} לא ריקה( ונוכיח ש-\L{$\bigcap\mathcal{C}$}
לא ריקה על ידי זה שנציג איבר קונקרטי \L{$x\in\bigcap\mathcal{C}$}.

בשלב ראשון, נרחיב את \L{$\mathcal{C}$} אל \L{$\mathcal{C}^{\prime}$}
מקסימלית שמקיימת את תכונת החיתוכים הסופיים. מכיוון ש-\L{$\mathcal{C}\subseteq\mathcal{C}^{\prime}$}
אז \L{$\bigcap\mathcal{C}^{\prime}\subseteq\bigcap\mathcal{C}$} ולכן
אם נוכיח ש-\L{$x\in\bigcap\mathcal{C}^{\prime}$} סיימנו; אז אפשר
פשוט להוכיח את המשפט מלכתחילה עבור \L{$\mathcal{C}$} מקסימלית ולא
לסחוב את התיוג הזה כל הזמן.

לכל \L{$A_{i}$} במכפלה שמגדירה את \L{$X$}, נסמן ב-\L{$\pi_{i}$}
את ההטלה \L{$\pi_{i}:X\to A_{i}$}. עכשיו נגדיר אוסף של קבוצות סגורות
ב-\L{$A_{i}$}: \L{$\mathcal{D}_{i}=\left\{ \overline{\pi_{i}\left(C\right)}\ |\ C\in\mathcal{C}\right\} $}.
כלומר: אנחנו לוקחים את הקבוצות של \L{$\mathcal{C}$}, מטילים אותן
אל \L{$A_{i}$}, ואז לוקחים את הסגור שלהן. עכשיו, \L{$\mathcal{D}_{i}$}
מקיימת גם היא את תכונת החיתוכים הסופיים: חיתוך של מספר סופי של איברים
ב-\L{$\mathcal{D}_{i}$} הוא מהצורה

\L{$\overline{\pi_{i}\left(C_{1}\right)}\cap\ldots\cap\overline{\pi_{i}\left(C_{n}\right)}$}

עבור \L{$C_{1},\ldots,C_{n}\in\mathcal{C}$}. ההנחה שאיתה התחלנו הייתה
ש-\L{$\mathcal{C}$} מקיימת את תכונת החיתוכים הסופיים, ולכן יש \L{$c\in C_{1}\cap\ldots\cap C_{n}$}
ומכאן ש-\L{$\pi_{i}\left(c\right)\in\overline{\pi_{i}\left(C_{1}\right)}\cap\ldots\cap\overline{\pi_{i}\left(C_{n}\right)}$}
)צריך לזכור שלקחת סגור של קבוצה רק מגדיל אותה(.

מכיוון ש-\L{$\mathcal{D}_{i}$} מקיימת את תכונת החיתוכים הסופיים \textbf{וגם}
\L{$A_{i}$} הוא מרחב קומפקטי, קיים \L{$d_{i}\in\bigcap\mathcal{D}_{i}$}.
את הדבר הזה נעשה לכל \L{$A_{i}$} שמשתתף במכפלה שמגדירה את \L{$X$},
ולסיום נגדיר \L{$x=\left(d_{i}\right)_{i}$}, כלומר \L{$x$} הוא האיבר
שבקואורדינטה ה-\L{$i$}-ית שלו יש את \L{$d_{i}$}.

כל מה שנשאר לעשות הוא להוכיח ש-\L{$x\in\bigcap\mathcal{C}$}, כלומר
ש-\L{$x\in C$} לכל \L{$C\in\mathcal{C}$}.

איך עושים את זה? מאיפה להתחיל?

שימו לב שעד כה בכלל לא השתמשתי בזה ש-\L{$C$} \textbf{סגורה}; כשהסתכלתי
על \L{$\pi_{i}\left(C\right)$} קודם לקחתי את הסגור, ואת זה אפשר לעשות
לכל קבוצה. האופן שבו אני אשתמש בזה הוא על ידי כך שלא אוכיח ישירות
ש-\L{$x\in C$} אלא אוכיח ש-\L{$x\in\overline{C}$} ואשתמש בכך ש-\L{$C$}
סגורה ולכן \L{$\overline{C}=C$}. כדי להוכיח ש-\L{$x\in\overline{C}$}
מספיק להראות שלכל סביבה של \L{$x$} יש חיתוך לא ריק עם \L{$C$}; ולא
באמת צריך להראות את זה לכל סביבה, מספיק להראות את זה לאברי הבסיס של
מרחב המכפלה \L{$X$}. בשביל להראות את זה מספיק להראות שכל איבר בסיס
של \L{$X$} שמכיל את \L{$x$} \textbf{שייך} אל \L{$\mathcal{C}$}
בעצמו. למה? כי אז, לכל \L{$C$} שניקח, החיתוך של \L{$C$} עם איבר
הבסיס הזה הוא לא ריק כי \L{$\mathcal{C}$} מקיימת את \textbf{תכונת
החיתוכים הסופיים}.

עכשיו אנחנו חוזרים אל שתי הטענות הכלליות שהוכחתי קודם ונעזרו בכך ש-\L{$\mathcal{C}$}
מקסימלית:
\begin{itemize}
\item אם \L{$\mathcal{C}$} מקסימלית ביחס לתכונת החיתוכים הסופיים אז היא
סגורה תחת חיתוכים סופיים.
\item אם \L{$\mathcal{C}$} מקסימלית ביחס לתכונת החיתוכים הסופיים ואם \L{$D$}
היא קבוצה שיש לה חיתוך לא ריק עם \textbf{כל} אבר ב-\L{$\mathcal{C}$},
אז \L{$D\in\mathcal{C}$}.
\end{itemize}
אנחנו רוצים להראות שכל איבר \textbf{בסיס} שמכיל את \L{$x$} שייך אל
\L{$\mathcal{C}$}, אבל בזכות הטענה הראשונה מספיק שנראה את זה לכל
איבר \textbf{תת-בסיס} שמכיל את \L{$x$}. כי כל איבר בסיס שמכיל את
\L{$x$} מתקבל מחיתוך של אברי תת-בסיס שכולם מכילים את \L{$x$}, ולכן
אם כולם היו שייכים אל \L{$\mathcal{C}$} כך גם איבר הבסיס.

עכשיו זה המקום שבו אנחנו משתמשים בכך שהטופולוגיה על \L{$X$} היא טופולוגיית
\textbf{המכפלה} ולא טופולוגיית \textbf{הקופסה}. בטופולוגיית המכפלה,
אברי תת-הבסיס הם פשוטים במיוחד: \L{$U_{i}=\pi^{-1}_{i}\left(V_{i}\right)$}
עבור כל קבוצה פתוחה \L{$V_{i}\subseteq A_{i}$}. כל מה שנשאר לנו לעשות
הוא להראות שאם \L{$x\in U_{i}$} אז \textbf{לכל} איבר ב-\L{$\mathcal{C}$}
יש חיתוך לא ריק עם \L{$U_{i}$}.

במבט ראשון, לא ברור לי איך עושים את זה - מה לקבוצה \L{$U_{i}$} ו\textbf{לכל}
אברי \L{$\mathcal{C}$}? צריך להיזכר מה קרה במהלך ההוכחה. מה שאנחנו
יודעים על \L{$U_{i}$} הוא ש-\L{$x$} שם; ו-\L{$x$} הוא לא סתם איבר,
הוא האיבר שבנינו בזהירות קואורדינטה-קואורדינטה, כשהקואורדינטה ה-\L{$i$}-ית
היא \L{$d_{i}\in\bigcap\mathcal{D}_{i}.$}, כאשר \L{$\mathcal{D}_{i}=\left\{ \overline{\pi_{i}\left(C\right)}\ |\ C\in\mathcal{C}\right\} $}-
הנה לנו הקשר של \L{$U_{i}$} אל \textbf{כל} הקבוצות ב-\L{$\mathcal{C}$}.

תהא אם כן \L{$C\in\mathcal{C}$} כלשהי ונסמן \L{$D_{i}=\pi_{i}\left(C\right)$}.
אז \L{$d_{i}\in\overline{D_{i}}$}. כעת, בנוסף לכך, \L{$d_{i}\in V_{i}$},
כלומר יש לנו סביבה \L{$V_{i}$} של \L{$d_{i}$}, ומההגדרה של סגור
של קבוצה ומכך ש-\L{$d_{i}\in\overline{D_{i}}$} נקבל שהחיתוך \L{$V_{i}\cap D_{i}$}
לא ריק, כלומר יש \L{$y\in C$} שמקיים \L{$\pi_{i}\left(y\right)\in V_{i}$},
ולכן \L{$y\in\pi^{-1}_{i}\left(V_{i}\right)=U_{i}$} וקיבלנו בדיוק
את מה שחיפשנו: \L{$y\in U_{i}\cap C$}, ולכן יש חיתוך לא ריק, כמבוקש.
זה מסיים באופן קצת מפתיע את כל ההוכחה של משפט טיכונוף.

בואו נסכם ונראה מה עשינו:
\begin{itemize}
\item בחרנו להוכיח את הקומפקטיות של \L{$X$} בעזרת ההגדרה שמתבססת על קבוצות
סגורות ותכונת החיתוך הסופי.
\item לקחנו אוסף \L{$\mathcal{C}$} של קבוצות שמקיים את תכונת החיתוך הסופי
והרחבנו אותו לאוסף \textbf{מקסימלי} שמקיים את התכונה הזו, בעזרת הלמה
של צורן.
\item בנינו \L{$x$} על ידי כך שהסתכלנו בכל ההטלות \L{$\pi_{i}$}, ולכל
הטלה כזו הפעלנו אותה על אברי \L{$\mathcal{C}$}, לקחנו את הסגור של
התוצאה, חתכנו את הכל וקיבלנו )בזכות הקומפקטיות של הרכיב ה-\L{$i$}(
איבר \L{$d_{i}$} שנמצא בחיתוך. \L{$x$} נבנה כך שבקואורדינטה ה-\L{$i$}
שלו יש את \L{$d_{i}$}.
\item הראינו שלכל אברי התת-בסיס הסטנדרטי של טופולוגיית המכפלה, אם הם מכילים
את \L{$x$} הם שייכים אל \L{$\mathcal{C}$}. זה נעזר במקסימליות של
\L{$\mathcal{C}$}.
\item אז השתמשנו במקסימליות של \L{$\mathcal{C}$} כדי לקבל מזה שכל איבר
בסיס של \L{$X$} שמכיל את \L{$x$} שייך אל \L{$\mathcal{C}$} בעצמו,
ולכן החיתוך שלו עם כל איבר ב-\L{$\mathcal{C}$} לא ריק שכן \L{$\mathcal{C}$}
מקיימת את תכונת החיתוכים הסופיים.
\item הסקנו שלכל \L{$C\in\mathcal{C}$} ולכל סביבה של \L{$x$} יש חיתוך
לא ריק עם \L{$C$}; לכן \L{$x\in\overline{C}$} ומכך ש-\L{$C$} סגורה
קיבלנו ש-\L{$x\in C$}, ומכאן ש-\L{$x\in\bigcap\mathcal{C}$}, כפי
שרצינו.
\end{itemize}
זו לא הוכחה שבאה לי טבעי מהראש )אני מסתבך יפה מאוד בחלקים האחרונים,
למרות שאת הקטע עם הלמה של צורן אני זוכר( אבל היא לא כזו מסובכת כפי
שהיה אפשר לחשוש!
\end{document}
